%% 自動生成（tools/mksteam2.py）。直接編集しないこと。値は IAPWS-IF97 の式（iapws 1.5.5）・CoolProp 8.0.0
\phbkhead{R32（ジフルオロメタン。家庭用エアコン）の飽和表}
{\footnotesize\color{bkGray}臨界点 \qty{78.1}{\celsius}・\qty{5.78}{MPa}。\par}\vspace{0.5mm}
\begin{longtblr}{colspec={Q[c,m,wd=9mm] Q[r,m] Q[r,m] Q[r,m] Q[r,m] Q[r,m] Q[r,m] Q[r,m]}, width=\linewidth, rowhead=1, colsep=3.2pt, rowsep=0.35pt, hlines={0.3pt,bkLine}, vline{2}={0.4pt,bkLine}, row{1-1}={bg=bkVol!14, font=\sffamily\bfseries\scriptsize, halign=c}, cells={font=\scriptsize}}
{温度\\$\theta$〔℃〕} & {圧力\\$p$〔MPa〕} & {液の密度\\$\rho'$〔kg/m$^3$〕} & {$v''$\\〔m$^3$/kg〕} & {$h'$\\〔kJ/kg〕} & {$h''$\\〔kJ/kg〕} & {$s'$\\〔kJ/(kg·K)〕} & {$s''$\\〔kJ/(kg·K)〕} \\
$-$40 & 0.1774 & 1180.2 & 0.1974 & 133.2 & 502.0 & 0.7382 & 2.3200 \\
$-$35 & 0.2214 & 1165.7 & 0.1601 & 141.3 & 504.2 & 0.7723 & 2.2962 \\
$-$30 & 0.2734 & 1151.0 & 0.1309 & 149.4 & 506.3 & 0.8060 & 2.2735 \\
$-$25 & 0.3346 & 1135.9 & 0.1079 & 157.7 & 508.2 & 0.8392 & 2.2518 \\
$-$20 & 0.4058 & 1120.6 & 0.08963 & 165.9 & 510.0 & 0.8720 & 2.2310 \\
$-$15 & 0.4881 & 1104.9 & 0.07493 & 174.3 & 511.6 & 0.9044 & 2.2109 \\
$-$10 & 0.5826 & 1088.8 & 0.06301 & 182.8 & 513.0 & 0.9365 & 2.1915 \\
$-$5 & 0.6906 & 1072.2 & 0.05328 & 191.3 & 514.3 & 0.9684 & 2.1727 \\
0 & 0.8131 & 1055.3 & 0.04527 & 200.0 & 515.3 & 1.0000 & 2.1543 \\
5 & 0.9514 & 1037.7 & 0.03862 & 208.8 & 516.1 & 1.0314 & 2.1363 \\
10 & 1.107 & 1019.7 & 0.03308 & 217.7 & 516.7 & 1.0628 & 2.1185 \\
15 & 1.281 & 1000.9 & 0.02842 & 226.8 & 516.9 & 1.0940 & 2.1008 \\
20 & 1.475 & 981.4 & 0.02448 & 236.1 & 516.9 & 1.1253 & 2.0831 \\
25 & 1.690 & 961.0 & 0.02112 & 245.6 & 516.5 & 1.1566 & 2.0652 \\
30 & 1.928 & 939.6 & 0.01826 & 255.3 & 515.7 & 1.1881 & 2.0471 \\
35 & 2.190 & 917.0 & 0.01579 & 265.3 & 514.5 & 1.2198 & 2.0285 \\
40 & 2.478 & 893.0 & 0.01365 & 275.6 & 512.7 & 1.2520 & 2.0091 \\
45 & 2.795 & 867.3 & 0.01178 & 286.3 & 510.3 & 1.2847 & 1.9888 \\
50 & 3.141 & 839.3 & 0.01015 & 297.5 & 507.1 & 1.3183 & 1.9670 \\
55 & 3.520 & 808.3 & 0.008697 & 309.3 & 502.9 & 1.3531 & 1.9432 \\
60 & 3.933 & 773.3 & 0.007396 & 321.9 & 497.4 & 1.3898 & 1.9166 \\
\end{longtblr}
\phbkhead{R134a（テトラフルオロエタン。カーエアコン・冷蔵庫）の飽和表}
{\footnotesize\color{bkGray}臨界点 \qty{101.1}{\celsius}・\qty{4.06}{MPa}。\par}\vspace{0.5mm}
\begin{longtblr}{colspec={Q[c,m,wd=9mm] Q[r,m] Q[r,m] Q[r,m] Q[r,m] Q[r,m] Q[r,m] Q[r,m]}, width=\linewidth, rowhead=1, colsep=3.2pt, rowsep=0.35pt, hlines={0.3pt,bkLine}, vline{2}={0.4pt,bkLine}, row{1-1}={bg=bkVol!14, font=\sffamily\bfseries\scriptsize, halign=c}, cells={font=\scriptsize}}
{温度\\$\theta$〔℃〕} & {圧力\\$p$〔MPa〕} & {液の密度\\$\rho'$〔kg/m$^3$〕} & {$v''$\\〔m$^3$/kg〕} & {$h'$\\〔kJ/kg〕} & {$h''$\\〔kJ/kg〕} & {$s'$\\〔kJ/(kg·K)〕} & {$s''$\\〔kJ/(kg·K)〕} \\
$-$40 & 0.05121 & 1417.7 & 0.3611 & 148.1 & 374.0 & 0.7956 & 1.7643 \\
$-$35 & 0.06614 & 1403.1 & 0.2840 & 154.4 & 377.2 & 0.8223 & 1.7575 \\
$-$30 & 0.08438 & 1388.4 & 0.2259 & 160.8 & 380.3 & 0.8486 & 1.7515 \\
$-$25 & 0.1064 & 1373.4 & 0.1816 & 167.2 & 383.4 & 0.8746 & 1.7461 \\
$-$20 & 0.1327 & 1358.3 & 0.1474 & 173.6 & 386.6 & 0.9002 & 1.7413 \\
$-$15 & 0.1639 & 1342.8 & 0.1207 & 180.1 & 389.6 & 0.9256 & 1.7371 \\
$-$10 & 0.2006 & 1327.1 & 0.09959 & 186.7 & 392.7 & 0.9506 & 1.7334 \\
$-$5 & 0.2433 & 1311.1 & 0.08280 & 193.3 & 395.7 & 0.9754 & 1.7300 \\
0 & 0.2928 & 1294.8 & 0.06931 & 200.0 & 398.6 & 1.0000 & 1.7271 \\
5 & 0.3497 & 1278.1 & 0.05837 & 206.8 & 401.5 & 1.0243 & 1.7245 \\
10 & 0.4146 & 1261.0 & 0.04944 & 213.6 & 404.3 & 1.0485 & 1.7221 \\
15 & 0.4884 & 1243.4 & 0.04209 & 220.5 & 407.1 & 1.0724 & 1.7200 \\
20 & 0.5717 & 1225.3 & 0.03600 & 227.5 & 409.7 & 1.0962 & 1.7180 \\
25 & 0.6654 & 1206.7 & 0.03091 & 234.5 & 412.3 & 1.1199 & 1.7162 \\
30 & 0.7702 & 1187.5 & 0.02664 & 241.7 & 414.8 & 1.1435 & 1.7145 \\
35 & 0.8870 & 1167.5 & 0.02303 & 249.0 & 417.2 & 1.1670 & 1.7128 \\
40 & 1.017 & 1146.7 & 0.01997 & 256.4 & 419.4 & 1.1905 & 1.7111 \\
45 & 1.160 & 1125.1 & 0.01734 & 263.9 & 421.5 & 1.2139 & 1.7092 \\
50 & 1.318 & 1102.3 & 0.01509 & 271.6 & 423.4 & 1.2375 & 1.7072 \\
55 & 1.492 & 1078.3 & 0.01314 & 279.5 & 425.2 & 1.2611 & 1.7050 \\
60 & 1.682 & 1052.9 & 0.01144 & 287.5 & 426.6 & 1.2848 & 1.7024 \\
\end{longtblr}
\phbkhead{R600a（イソブタン。家庭用冷蔵庫）の飽和表}
{\footnotesize\color{bkGray}臨界点 \qty{134.7}{\celsius}・\qty{3.63}{MPa}。\par}\vspace{0.5mm}
\begin{longtblr}{colspec={Q[c,m,wd=9mm] Q[r,m] Q[r,m] Q[r,m] Q[r,m] Q[r,m] Q[r,m] Q[r,m]}, width=\linewidth, rowhead=1, colsep=3.2pt, rowsep=0.35pt, hlines={0.3pt,bkLine}, vline{2}={0.4pt,bkLine}, row{1-1}={bg=bkVol!14, font=\sffamily\bfseries\scriptsize, halign=c}, cells={font=\scriptsize}}
{温度\\$\theta$〔℃〕} & {圧力\\$p$〔MPa〕} & {液の密度\\$\rho'$〔kg/m$^3$〕} & {$v''$\\〔m$^3$/kg〕} & {$h'$\\〔kJ/kg〕} & {$h''$\\〔kJ/kg〕} & {$s'$\\〔kJ/(kg·K)〕} & {$s''$\\〔kJ/(kg·K)〕} \\
$-$40 & 0.02870 & 624.1 & 1.143 & 112.5 & 501.3 & 0.6549 & 2.3227 \\
$-$35 & 0.03680 & 618.9 & 0.9074 & 123.0 & 507.9 & 0.6995 & 2.3154 \\
$-$30 & 0.04662 & 613.6 & 0.7284 & 133.7 & 514.4 & 0.7437 & 2.3095 \\
$-$25 & 0.05843 & 608.3 & 0.5906 & 144.4 & 521.0 & 0.7874 & 2.3048 \\
$-$20 & 0.07248 & 602.9 & 0.4834 & 155.3 & 527.6 & 0.8306 & 2.3013 \\
$-$15 & 0.08905 & 597.4 & 0.3990 & 166.3 & 534.3 & 0.8735 & 2.2989 \\
$-$10 & 0.1085 & 591.9 & 0.3320 & 177.4 & 540.9 & 0.9160 & 2.2975 \\
$-$5 & 0.1310 & 586.3 & 0.2783 & 188.6 & 547.6 & 0.9582 & 2.2969 \\
0 & 0.1570 & 580.6 & 0.2349 & 200.0 & 554.3 & 1.0000 & 2.2972 \\
5 & 0.1867 & 574.8 & 0.1995 & 211.5 & 561.1 & 1.0415 & 2.2983 \\
10 & 0.2206 & 568.9 & 0.1704 & 223.1 & 567.8 & 1.0828 & 2.3000 \\
15 & 0.2590 & 562.9 & 0.1464 & 234.9 & 574.5 & 1.1239 & 2.3023 \\
20 & 0.3022 & 556.9 & 0.1264 & 246.9 & 581.2 & 1.1647 & 2.3051 \\
25 & 0.3507 & 550.7 & 0.1096 & 259.0 & 587.9 & 1.2053 & 2.3085 \\
30 & 0.4047 & 544.3 & 0.09542 & 271.2 & 594.6 & 1.2458 & 2.3123 \\
35 & 0.4648 & 537.8 & 0.08341 & 283.7 & 601.2 & 1.2861 & 2.3165 \\
40 & 0.5312 & 531.2 & 0.07317 & 296.3 & 607.8 & 1.3263 & 2.3211 \\
45 & 0.6044 & 524.4 & 0.06440 & 309.1 & 614.3 & 1.3664 & 2.3259 \\
50 & 0.6849 & 517.4 & 0.05683 & 322.1 & 620.8 & 1.4064 & 2.3309 \\
55 & 0.7730 & 510.2 & 0.05029 & 335.3 & 627.2 & 1.4464 & 2.3361 \\
60 & 0.8692 & 502.7 & 0.04459 & 348.7 & 633.5 & 1.4863 & 2.3414 \\
\end{longtblr}
